\documentclass[12pt]{article}
\usepackage[margin=0.75in]{geometry}
\geometry{a4paper}
\usepackage[T1]{fontenc} 
\usepackage[utf8]{inputenc} 
\usepackage{graphicx} 
\usepackage{hyperref} 
\usepackage{amsmath}
\usepackage{multicol}
\usepackage{tcolorbox}
\usepackage{xcolor}
\usepackage{tikz}

% Define a custom command for fill-in-the-blank
\newcommand{\fillin}[2][4cm]{%
  \begin{tikzpicture}[baseline=(current bounding box.south)]
    \draw[thick] (0,0) -- (#1,0);
    \if\relax\detokenize{#2}\relax
      % Do nothing if #2 is empty
    \else
      \node[above=2pt] at (0.5*#1,0) {#2};
    \fi
  \end{tikzpicture}%
}


\title{
\centerline{\includegraphics[width=50mm]{images/ru_logo.jpg}}
\vspace{0.5 cm}
Einstaklingsverkefni 2 (e. Individual Assignment 2)
\large  \\
T-117-STR1, Strjál stærðfræði I, 2024-3 \\ 
\small Reykjavik University - Department of Computer Science, Menntavegi 1, IS-101 Reykjavík, Iceland 
  }

\author{
    Kennari: Harpa Guðjónsdóttir\\
    \texttt{harpagud@ru.is}
}

\date{Skilafrestur (e. Deadline): 10.09.2024}

\begin{document}
\maketitle
\noindent
Hér er Einstaklingsverkefni 2. Skilafrestur er þriðjudaginn 10.september 2024 kl. 23:59*. Þetta eru ein af 5 einstaklingsskilum. Þau gilda alls 20\% af lokaeinkunn, en lægstu einkunn er sleppt.
\vspace{2mm}
\newline
Mjög mikilvægt er að nemendur noti þetta skjal, fylli inn sínar lausnir á viðeigandi staði og skili útfylltu skjali á Gradescope sem pdf. Bæði er leyfilegt að prenta út skjalið, fylla inn handvirkt og skanna það svo aftur inn (\textbf{eða} nota þetta \LaTeX{} sniðmát og fylla inn í það). \textbf{Ekki verður farið yfir verkefni sem ekki nota þetta skjal (eða \LaTeX{} sniðmátið), og fyrir slík verkefni fæst 0 í einkunn.}
\vspace{2mm}
\newline
\footnotesize
*nemendur á Austurlandi skila miðvikudaginn 11.september 2023 kl.23:59 og skilafrestur þeirra í Canvas/Gradescope er stilltur miðað við það
\normalsize	
\subsubsection*{English version:}
("This is the second individual assignment. The deadline is Tuesday, September 10th, 2024, at 23:59*. Students hand in solutions on pdf on Gradescope. This is one of 5 individual assignments. All in all, their weight is 20\% of the final grade, but the lowest grade is dropped.")
\vspace{2mm}
\newline
Students must use this document, fill in their solutions in the designated spaces, and return the completed document to Gradescope as a pdf. You are allowed to print the document, fill it in writing, and scan it, (\textbf{or} use this \LaTeX{} template and fill it in). \textbf{Assignments solutions that do not use this document (or the \LaTeX{} template) will not be reviewed and will receive a grade of 0.}
\vspace{2mm}
\newline
\footnotesize
*Students in the east of Iceland hand in on Wednesday, September 11th, 2024, at 23:59, and their deadline is set in Canvas/Gradescope accordingly
\normalsize	
\newpage
\section*{\textbf{Skiladæmi (e. Hand-in problems) :}}
\subsection*{Dæmi 1 (e. Problem 1) (14\%)\label{section:problem1}}
Búið til sanntöflu fyrir samsettu yrðinguna $\neg(r\land \neg q) \to(p\lor \neg r)$. Sýnið öll milliskref í töflunni, og munið að fylla inn í hausinn á töflunni. Er yrðingin sísanna? \\ \\
("Write the truth table for the proposition $\neg(r\land \neg q) \to(p\lor \neg r)$. Show all intermediate columns, and remember to fill in the head row. Is the proposition a tautology?")

\subsection*{Svar við Dæmi 1 (e. Answer to Problem 1)}

\begin{center}
    \begin{tabular}{ |c|c|c|p{0.8 cm}|p{0.8 cm}|p{1.8 cm}|p{2 cm}|p{2 cm}|p{3.8 cm}|} 
        \hline
        $p$ & $q$ & $r$ & & & & & & \\
        \hline
        T & T & T &  &  &  &  &  & \\ 
        \hline
        T & T & F &  &  &  &  &  & \\ 
        \hline
        T & F & T &  &  &  &  &  & \\ 
        \hline
        T & F & F &  &  &  &  &  & \\ 
        \hline
        F & T & T &  &  &  &  &  & \\
        \hline
        F & T & F &  &  &  &  &  & \\ 
        \hline
        F & F & T &  &  &  &  &  & \\ 
        \hline
        F & F & F &  &  &  &  &  & \\
        \hline
    \end{tabular}
\end{center}
Er yrðingin sísanna? ("Is the proposition a tautology?")
\fillin[5cm]{
    % Start your answer to a) after this line:
  
    % End your answer to a) before this line.
    }     
\subsection*{Dæmi 2 (e. Problem 2) (20\%)\label{section:daemi2}}
Notið \textbf{sönnun með mótskilyrðingu} til að sanna eftirfarandi regluna. \\
("Use \textbf{proof by contraposition} to prove the following rule.")\\ \\
\indent Regla: Ef $n$ er heiltala og $5n - 1$ er slétt tala þá er $n$ oddatala. \\
\indent ("Rule: If $n$ is an integer and $5n-1$ is even, then $n$ is odd.")

\subsection*{Svar við Dæmi 2 (e. Answer to Problem 2)}

\newcommand{\makeAnswerBoxProblemTwo}{
  \noindent
  \begin{tcolorbox}[colframe=black, colback=white, boxrule=0.5pt, arc=0pt, outer arc=0pt, height=\dimexpr\textheight-\ht\strutbox-9\baselineskip\relax]
    % Start your answer after this line:

    
    % End your answer before this line.
  \end{tcolorbox}
}

\makeAnswerBoxProblemTwo



\newpage
\subsection*{Dæmi 3 (e. Problem 3) (8\%+8\%)\label{section:daemi3}}
\textbf{Skilgreining á NAND og NOR:} Yrðingin p NAND q er sönn þegar annað hvort p eða q, eða báðar, eru ósannar; og yrðingin er ósönn þegar bæði p og q eru sannar. Yrðingin p NOR q er sönn þegar bæði p og q eru ósannar, og yrðingin er ósönn að öðru leyti. Yrðingarnar p NAND q og p NOR q eru táknaðir með $p | q$ og $p \downarrow q$, í þessari röð. \\ \\
("\textbf{Definition of NAND and NOR:} The proposition p NAND q is true when either p or q, or both, are false; and it is false when both p and q are true. The proposition p NOR q is true when both p and q are false, and it is false otherwise. The propositions p NAND q and p NOR q are denoted by $p | q$ and $p \downarrow q$, respectively.")

\begin{itemize}
    \item[a)] Setjið upp sanntöflu fyrir rökvirkjann NAND. ("Construct a truth table for the logical operator NAND.")
    \item[b)] Setjið upp sanntöflu fyrir rökvirkjann NOR. ("Construct a truth table for the logical operator NOR.")
\end{itemize}

\subsection*{Svör við Dæmi 3 (e. Answers to Problem 3)}

\begin{itemize}
    \item[a)]
        \begin{minipage}[t]{\linewidth}
            \renewcommand{\arraystretch}{1.5} 
            \setlength{\tabcolsep}{12pt} 
            \begin{tabular}{ |c|c|p{1.5 cm}| }
                \hline
                $p$ & $q$ & $p | q$\\
                \hline
                T & T &  \\
                \hline
                T & F &  \\
                \hline
                F & T &  \\
                \hline
                F & F &  \\
                \hline
            \end{tabular}
        \end{minipage}
    \item[b)]
        \begin{minipage}[t]{\linewidth}
            \renewcommand{\arraystretch}{1.5} 
            \setlength{\tabcolsep}{12pt} 
            \begin{tabular}{ |c|c|p{1.5 cm}| }
                \hline
                $p$ & $q$ & $p \downarrow q$\\
                \hline
                T & T &  \\
                \hline
                T & F &  \\
                \hline
                F & T &  \\
                \hline
                F & F &  \\
                \hline
            \end{tabular}
        \end{minipage}
\end{itemize}
\newpage
\subsection*{Dæmi 4 (e. Problem 4) (14\%+10\%) \label{section:daemi4}}
Könnun meðal 500 sjónvarpsáhorfenda gefur eftirfarandi upplýsingar. Af þeim horfa 285 á fótbolta, 195 á handbolta og 115 á körfubolta. Það horfa 70 á fótbolta og handbolta, 45 á fótbolta og körfubolta og 42 á handbolta og körfubolta. Það eru 50 sem horfa ekki á neina af þessum íþróttagreinum. \\
("A survey of 500 television viewers provides the following information. Of them 285 watch soccer, 195 watch handball and 115 watch basketball. Also, 70 watch soccer and handball, 45 watch soccer and basketball and 42 watch in handball and basketball. There are 50 viewers who do not watch any of these sports.")
\begin{itemize}
    \item[a)] Hve margir þátttakendur í könnuninni horfa á allar íþróttagreinarnar? ("ow manu watch all three sports?")
    \item[b)] Hversu margir horfa á fótbolta en hvorki handbolta né körfubolta? ("How many watch soccer but neither handball nor basketball?")
\end{itemize}
Sýnið alla útreikninga í svarkassanum, og setjið að auki lokasvörin á línurnar fyrir neðan. ("Show all calculations in the answer box, and additionally put your final answers on the lines below.")

\subsection*{Svör við Dæmi 4 (e. Answers to Problem 4)}

\newcommand{\makeAnswerBoxProblemFour}{
  \noindent
  \begin{tcolorbox}[colframe=black, colback=white, boxrule=0.5pt, arc=0pt, outer arc=0pt, height=\dimexpr\textheight-\ht\strutbox-25\baselineskip\relax]
    % Start your answer after this line:

    
    % End your answer before this line.
  \end{tcolorbox}
}

\makeAnswerBoxProblemFour 

\vspace{5mm}

\begin{multicols}{2}
\begin{itemize}
    \item [a)] \fillin[7cm]{
    % Start your final answer to a) after this line:
  
    % End your final answer to a) before this line.
    }     
    \newline
    \item [b)] \fillin[7cm]{
    % Start your final answer to b) after this line:
  
    % End your final answer to b) before this line.
    } 
    \newline   
\end{itemize}
\end{multicols}
\newpage
\subsection*{Dæmi 5 (e. Problem 5) (4\%+4\%+4\%+4\%+4\%) \label{section:daemi5}}
Látum $I(x)$ vera yrðinguna “$x$ er með nettengingu” og $C(x, y)$ vera yrðinguna “$x$ og $y$ hafa talað saman á netinu,” þar sem mengi breytanna $x$ og $y$ eru allir nemendur í bekknum þínum. Ritið eftirfarandi staðhæfingar með því að nota eftir því sem við á allsherjarkvantara, tilvistarkvantara, rökfræðileg tákn og föllin $I(x)$ og $C(x,y)$\\
\\
("Let $I(x)$ be the statement “$x$ has an Internet connection” and $C(x, y)$ be the statement “$x$ and $y$ have chatted over the Internet,” where the domain for the variables $x$ and $y$ consists of all students in your class. Write the statements using these predicates, universal quantification, existential quantification, logical operators as well a $I(x)$ and $C(x,y)$.

\begin{itemize}
    \item[a)] Jón er ekki með nettengingu. ("Jón does not have an Internet connection.")
    \item[b)] Rúna hefur ekki talað við Sessu á netinu. ("Rúna has not chatted over the Internet with Sessa.")
    \item[c)] Enginn í bekknum hefur talað við Bóas á netinu. ("No one in the class has chatted with Bóas over the internet.")
    \item[d)] Allir nema einn nemandi í bekknum þínum er með nettengingu. ("Everyone except one student in your class has an Internet connection.")
    \item[e)] Það eru að minnsta kosti tveir nemendur í bekknum sem hafa ekki talað við sömu manneskjuna í bekknum á netinu.
    ("There are at least two students in your class who have not chatted over the internet with the same person in your class.")
\end{itemize}

\subsection*{Svör við Dæmi 5 (e. Answers to Problem 5)}

\newcommand{\makeAnswerBoxProblemFiveA}{
  \begin{tcolorbox}[colframe=black, colback=white, boxrule=0.5pt, arc=0pt, outer arc=0pt, height=1.8cm, width=\linewidth, top=1mm, bottom=1mm, left=1mm, right=1mm]
    % Start your answer to a) after this line:

    % End your answer to a) before this line.
  \end{tcolorbox}
}

\newcommand{\makeAnswerBoxProblemFiveB}{
  \begin{tcolorbox}[colframe=black, colback=white, boxrule=0.5pt, arc=0pt, outer arc=0pt, height=1.8cm, width=\linewidth, top=1mm, bottom=1mm, left=1mm, right=1mm]
    % Start your answer to b) after this line:

    % End your answer to b) before this line.
  \end{tcolorbox}
}

\newcommand{\makeAnswerBoxProblemFiveC}{
  \begin{tcolorbox}[colframe=black, colback=white, boxrule=0.5pt, arc=0pt, outer arc=0pt, height=1.8cm, width=\linewidth, top=1mm, bottom=1mm, left=1mm, right=1mm]
    % Start your answer to c) after this line:

    % End your answer to c) before this line.
  \end{tcolorbox}
}

\newcommand{\makeAnswerBoxProblemFiveD}{
  \begin{tcolorbox}[colframe=black, colback=white, boxrule=0.5pt, arc=0pt, outer arc=0pt, height=1.8cm, width=\linewidth, top=1mm, bottom=1mm, left=1mm, right=1mm]
    % Start your answer to d) after this line:

    % End your answer to d) before this line.
  \end{tcolorbox}
}

\newcommand{\makeAnswerBoxProblemFiveE}{
  \begin{tcolorbox}[colframe=black, colback=white, boxrule=0.5pt, arc=0pt, outer arc=0pt, height=1.8cm, width=\linewidth, top=1mm, bottom=1mm, left=1mm, right=1mm]
    % Start your answer to e) after this line:

    % End your answer to e) before this line.
  \end{tcolorbox}
}

\begin{itemize}
    \item[a)] \makeAnswerBoxProblemFiveA
    \item[b)] \makeAnswerBoxProblemFiveB
    \item[c)] \makeAnswerBoxProblemFiveC
    \item[d)] \makeAnswerBoxProblemFiveD
    \item[e)] \makeAnswerBoxProblemFiveE
\end{itemize}
\end{document}
